\documentclass[11pt]{article}
\usepackage{amsmath, amssymb}
\usepackage{geometry}
\usepackage{listings}
\usepackage{xcolor}
\usepackage{hyperref}
\geometry{margin=2.3cm}

\lstset{
    basicstyle=\ttfamily\small,
    columns=fullflexible,
    keepspaces=true,
    breaklines=true,
    frame=single,
    framerule=0.3pt,
    rulecolor=\color{gray!50},
    backgroundcolor=\color{gray!5},
    aboveskip=4pt,
    belowskip=4pt,
}

\newcommand{\file}[1]{\texttt{#1}}
\newcommand{\code}[1]{\texttt{#1}}
\newcommand{\dd}{\mathrm{d}}
\newcommand{\Tr}{\mathrm{Tr}}
\newcommand{\calN}{\mathcal{N}}
\newcommand{\Var}{\operatorname{Var}}
\newcommand{\tauint}{\tau_{\mathrm{int}}}

\title{Error bars for the pCN sampler: batch means (blocking) and the autocorrelation time}
\date{}

\begin{document}
\maketitle

\section{The problem}

The pCN sampler produces, per MPI core, a Markov chain of mean-field trajectories
$V$ and from it a stream of single-sample observables
\[
  o_1, o_2, \dots, o_L ,
\qquad
  o_i = \frac{\Tr\!\big[U(\beta-\tau;V_i)\,S^a\,U(\tau;V_i)\,S^b\big]}{Z(V_i)} ,
\]
i.e.\ one such stream for every correlation point $(a,b,\tau_k)$ (real and
imaginary part). The reported correlation is the global mean over all
$N = C\,L$ samples ($C$ = number of cores, each an independent chain),
\[
  \bar o = \frac1N \sum_{i=1}^{N} o_i .
\]
We need the standard error of $\bar o$.

If the samples were independent, the squared standard error would be
\begin{equation}
  \sigma_{\mathrm{iid}}^2 = \frac{s^2}{N},
  \qquad
  s^2 = \frac1{N}\sum_i o_i^2 - \bar o^{\,2}
  \label{eq:iid}
\end{equation}
with $s^2$ the marginal sample variance. But a Markov chain produces
\emph{autocorrelated} samples. Writing $\rho_k$ for the normalised
autocorrelation of $o$ at lag $k$ and
\begin{equation}
  \tauint = \tfrac12 + \sum_{k\ge 1}\rho_k
  \label{eq:tauint-def}
\end{equation}
for the integrated autocorrelation time (in units of samples; $\tauint=\tfrac12$
for independent draws), the true variance of the mean is inflated:
\begin{equation}
  \Var(\bar o) = \frac{s^2}{N}\,\big(2\tauint\big)
  = \sigma_{\mathrm{iid}}^2 \cdot 2\tauint .
  \label{eq:inflation}
\end{equation}
Equivalently the effective sample size is $N_{\mathrm{eff}} = N/(2\tauint)$.
Using \eqref{eq:iid} instead of \eqref{eq:inflation} underestimates the error by
$\sqrt{2\tauint}$, which at low temperature is large (here $\tauint\sim10^3$,
Sec.~\ref{sec:numbers}). The previous estimator scaled $\sigma_{\mathrm{iid}}$ by
an AR(1) factor inferred from the acceptance rate; that captures only the lag-1
term $\rho_1$ and badly underestimates $\tauint$ when the chain has a slow mode.
The fix is to measure $\Var(\bar o)$ directly, without modelling $\rho_k$.

\section{Batch means (blocking)}

Split each chain's $L$ samples into $B$ consecutive blocks of length
$\ell = \lfloor L/B\rfloor$ and form the block means
\begin{equation}
  \bar o_b = \frac1\ell \sum_{i\in \text{block }b} o_i ,
  \qquad b = 1,\dots,B .
\end{equation}
Adjacent samples within a block are correlated, but two block means become
\emph{independent} once the block length exceeds the autocorrelation time,
$\ell \gg \tauint$. Pooling the block means of all cores gives
$M = C\,B$ approximately independent estimates of the same quantity, and the
standard error of their average --- which is exactly the global mean $\bar o$,
since the blocks tile the data with equal weight --- is
\begin{equation}
  \boxed{\;
  \hat\sigma_{\mathrm{block}}^2
  = \frac1M \,\widehat{\Var}\big(\{\bar o_m\}_{m=1}^{M}\big)
  = \frac1M\cdot\frac1{M-1}\sum_{m=1}^{M}\big(\bar o_m - \bar o\big)^2 .
  \;}
  \label{eq:batchmeans}
\end{equation}
This is the batch-means estimator. It needs no model of the autocorrelation: the
block length absorbs it. Crucially the statistics come from the \emph{number of
blocks} $M=CB$, not from the number of cores $C$, so it is usable with few cores
(here $4$--$16$ cores $\times\,32$ blocks $=128$--$512$ pooled block means).

\paragraph{Pooling.} Each core flattens its $B$ block means; an
\code{MPI\_Allgather} collects all $CB$ of them on every rank, which then evaluate
\eqref{eq:batchmeans} per correlation point.

\section{Validity check: the blocking curve and the plateau}
\label{sec:plateau}

Equation \eqref{eq:batchmeans} is only correct when $\ell \gg \tauint$; for too-short
blocks it \emph{underestimates} the error. The Flyvbjerg--Petersen test makes this
visible: recompute $\hat\sigma_{\mathrm{block}}$ for a sequence of block lengths by
repeatedly merging adjacent blocks \emph{within each chain},
$\ell \to 2\ell \to 4\ell \to \cdots$ (which keeps the merged blocks contiguous in
chain time). The estimate rises with $\ell$ and then \emph{plateaus} once
$\ell \gg \tauint$:
\[
  \ell \ll \tauint:\ \hat\sigma_{\mathrm{block}}\ \text{too small}
  \qquad\longrightarrow\qquad
  \ell \gg \tauint:\ \hat\sigma_{\mathrm{block}}\ \to\ \text{true } \sigma(\bar o).
\]
The code therefore does \emph{not} report the finest binning. For each point it
takes the \textbf{maximum} of $\hat\sigma_{\mathrm{block}}$ over all merge levels
whose pooled block count stays $\ge 8$ (the level $\ell=L/B$ is always included).
The curve is monotone-then-noisy, so the maximum is a robust read of the plateau,
and it never under-covers. The aggregate curve (mean over the computed
$\mathrm{Re}\,g$ points vs.\ block length) is written to the output so the plateau
can be inspected; the printed ``blocking plateau ratio'' (longest/shortest level)
is $\approx 1$ when the blocks are long enough and $\gg 1$ when the run is too
short for its autocorrelation.

\section{Recovering the autocorrelation time}

The blocking error and the i.i.d.\ error differ precisely by the inflation factor
\eqref{eq:inflation}. Once the blocking estimate has plateaued it equals the true
$\Var(\bar o)$, so the ratio gives $\tauint$ directly:
\begin{equation}
  2\tauint
  = \frac{\Var(\bar o)}{\sigma_{\mathrm{iid}}^2}
  = \frac{\hat\sigma_{\mathrm{block}}^2}{s^2/N}
  \qquad\Longrightarrow\qquad
  \boxed{\;
  \tauint = \frac12\,\frac{\hat\sigma_{\mathrm{block}}^2\,N}{s^2} .
  \;}
  \label{eq:tau-recover}
\end{equation}
Both ingredients are already in hand: $\hat\sigma_{\mathrm{block}}$ from
\eqref{eq:batchmeans} and the marginal variance $s^2$ from the running
$\sum_i o_i^2$ (the \code{sample\_sqsum} accumulator), so $\tauint$ costs nothing
extra. It is computed per correlation point.

\paragraph{Interpretation.} $\tauint$ from \eqref{eq:tau-recover} is in units of
pCN steps (samples) and is a property of a \emph{single} chain: because the cores
are independent, $N_{\mathrm{eff}} = C\,L/(2\tauint) = N/(2\tauint)$, and the
$C$-dependence cancels in \eqref{eq:tau-recover}. The independent limit gives
$\tauint=\tfrac12$; a large value means consecutive samples stay correlated over
$\sim\tauint$ steps. This is the number that sets the practical rules of thumb:
\begin{itemize}
  \item \textbf{burn-in} $\gtrsim 5$--$10\,\tauint$ (forget the initial condition);
  \item \textbf{per-core chain length} $L \gtrsim 10$--$20\,\tauint$ (so each chain
        is well mixed and unbiased, and \eqref{eq:batchmeans} has plateaued);
  \item under slow mixing prefer \emph{long} chains over \emph{many short} ones:
        initialization bias does not average across cores, and burn-in is paid per
        chain.
\end{itemize}

\section{In the code}

\begin{tabular}{@{}ll@{}}
\code{-{}-errmethod} & \code{blocking} (default, this note) or \code{ar1} (legacy acceptance factor)\\
\code{-{}-numBlocks} & number of blocks per core $B$ (default $32$)\\
\end{tabular}

\medskip
\noindent
Per-sample accumulation into the current block and block closing live in
\file{Run\_Time\_Data} (\code{init\_blocks}, \code{close\_block}) fed by
\code{compute\_spin\_correlations}; the pooling, plateau and $\tauint$ are in
\code{compute\_sample\_stds\_blocking}. Written to the HDF5 \code{runtimedata} group:
\begin{itemize}
  \item \code{Re/Im\_correlation\_sample\_stds} --- the plateau blocking error \eqref{eq:batchmeans};
  \item \code{Re/Im\_correlation\_tau\_int} --- $\tauint$ in steps \eqref{eq:tau-recover};
  \item \code{blocking\_curve\_block\_length}, \code{blocking\_curve\_sigma} --- the Flyvbjerg--Petersen curve;
  \item attribute \code{error\_method}.
\end{itemize}
Each iteration prints the ``blocking plateau ratio'' and ``max tau\_int (steps)''.

\section{Measured behaviour ($\beta=12$, ISO, symmetry type C)}
\label{sec:numbers}

The slow mode shows up in the should-be-zero $\mathrm{Re}\,g^{xy}(\tau)$. Validated
against the spread of independent runs (the ground-truth error):
\begin{itemize}
  \item $\tauint \approx 2800$ steps. This explains the sample-size dependence: at
        $L=10^4$/core the block length $L/B\approx310$ is $\ll\tauint$ and the error
        is underestimated; at $L=10^5$/core, $L/B\approx3100\gtrsim\tauint$ and the
        estimator is calibrated.
  \item At $L=10^5$/core the reported-to-true error ratio is $\approx 1.0$, and
        $\mathrm{Re}\,g^{xy}$ is statistically consistent with zero ($\lesssim 2\sigma$).
        The apparent ``symmetry violation'' seen with the AR(1) error was an
        underestimated error bar, not broken physics.
  \item The between-run spread falls $\sim 1/\sqrt N$ as samples grow, i.e.\ the chain
        is \emph{slow but ergodic}, not trapped behind a barrier --- more samples (and
        burn-in $\gtrsim5$--$10\,\tauint$) suffice; tempering is not required.
\end{itemize}

\end{document}
